\section{Introduction}
\label{sec:introduction}

\subsection{The DIA Administrator Question}

Suppose the \dia{} is enacted and the redistricting algorithm is run by
a federal administrator (a ``Districting Administrator'') with a fixed
parameter set encoded in statute.
A hostile legislature challenges the map in court, arguing that the
parameter set was chosen to produce a partisan outcome.
Alternatively, a well-meaning administrator asks: ``If I tweak the
METIS imbalance factor slightly, will the partisan outcome change?''

This paper answers both questions empirically through a comprehensive
50-state parameter sweep.
The central finding: \emph{partisan seat counts are insensitive to
parameter tuning}.
The maximum variation in Democratic seat count across the full parameter
range of all five parameters is $\Delta D\text{-seats} \leq 0.3$ nationally---
indistinguishable from sampling noise.

This finding complements B.0's bakeoff~\citep{b0paper}: while algorithm
\emph{structure} (e.g., GeoSection vs.\ unweighted bisection) can produce
partisan gaps of 12 percentage points, parameter \emph{tuning} within a
fixed algorithm produces negligible partisan effects.
The DIA statute can safely fix any reasonable parameter set.

\subsection{Parameters Studied}

We study five parameters that appear in the \ar{} + \cs{} pipeline:

\begin{enumerate}
\item \textbf{METIS imbalance factor} ($\ufactor \in [0.1\%, 5\%]$):
      the maximum allowable population deviation from perfect balance
      within the METIS heuristic.
\item \textbf{County edge weight} ($\acounty \in [1.0, 10.0]$):
      the multiplicative boost applied to edges that cross county lines.
\item \textbf{ConvergenceSweep threshold} ($T \in [100, 1{,}000]$):
      the consecutive non-improving seeds required to halt.
\item \textbf{Area swing} ($\aswing \in [1.05, 1.20]$):
      the area-to-population ratio parameter in GeoSection/AreaSection.
\item \textbf{VRA boost weight} ($\wvra \in [0.2, 0.6]$):
      the multiplicative factor for edges adjacent to VRA-eligible
      minority community tracts.
\end{enumerate}

The baseline parameter set is the DIA default: $\ufactor = 1\%$,
$\acounty = 3.0$, $T = 600$, $\aswing = 1.10$, $\wvra = 0.4$.

\subsection{Main Results}

\begin{enumerate}
\item \textbf{Partisan outcomes are insensitive.}
      Across all five parameters and their full ranges, the national
      Democratic seat count varies by $\leq 0.3$ seats.
      No parameter, alone or in combination, can systematically shift the
      partisan outcome by more than one seat nationally.
\item \textbf{Compactness is sensitive to $\ufactor$.}
      Polsby-Popper varies by $\pm 0.010$ across the $\ufactor$ range.
      Tighter $\ufactor$ (e.g., 0.1\%) produces more compact districts;
      looser $\ufactor$ (e.g., 5\%) allows more population imbalance
      but also more compactness.
\item \textbf{$T$ affects runtime, not outcomes.}
      The \cs{} threshold $T$ has no measurable effect on partisan
      outcomes across the range $[100, 1{,}000]$.
      This confirms B.16's finding that the algorithm converges well
      before $T = 600$.
\end{enumerate}

\subsection{Paper Structure}

Section~\ref{sec:parameters} defines the five parameters precisely.
Section~\ref{sec:methodology} describes the sweep design.
Section~\ref{sec:results} presents per-parameter results.
Section~\ref{sec:key-finding} synthesises the main finding.
Section~\ref{sec:implications} draws implications for the DIA statute.
Section~\ref{sec:conclusion} concludes.
